% ==========================================================
% 010_overview.tex
% Forecast Admissibility Surface (FAS)
% ==========================================================

\section{Overview}
\label{sec:fas_overview}

Operational forecasting pipelines often proceed under the implicit assumption
that all demand slices are amenable to learning, evaluation, and control once
sufficient data and model sophistication are applied. In practice, this
assumption frequently fails. Certain slices of demand—defined at the level of
entity, entity-by-interval, or site-by-entity-by-interval—exhibit structural
pathologies that render forecasting unstable or misleading regardless of model
class. These pathologies may arise from extreme intermittency, insufficient
support, heavy-tailed error distributions, or recurrent spike-driven behavior
that dominates baseline error anatomy.

When such slices are treated as fully admissible, downstream systems incur
several risks: models overfit noise, evaluation metrics become erratic,
governance signals fluctuate unpredictably, and readiness adjustments are
applied in regimes where learning is neither reliable nor actionable. In these
settings, additional modeling complexity does not resolve the issue; it merely
obscures the underlying structural hostility of the slice.

\medskip
\begin{framed}
\noindent\textbf{Core Framing Question.}\;
\emph{Given a particular slice of the forecast domain, is this slice structurally
admissible for learning, evaluation, and control, or does its observed error
anatomy indicate that forecasting should be constrained or avoided altogether?}
\end{framed}
\medskip

The \emph{Forecast Admissibility Surface (FAS)} is introduced as a deterministic,
slice-keyed construct that explicitly answers this question.

FAS is not a performance metric and is not an optimization objective. It does not
rank models, tune hyperparameters, or estimate forecast quality. Instead, it
classifies slices of the forecast domain into admissibility regions based on
observable properties of baseline error anatomy, including support size, spike
frequency, and tail magnitude. Its sole purpose is to determine \emph{where}
forecasting and readiness control may be meaningfully applied, and where such
efforts should be restricted or blocked entirely.

Within the Forecast Readiness Framework (FRF), FAS occupies a deliberately
upstream position. Whereas Demand Quantization Compatibility (DQC) assesses
whether realized demand admits stable interpretation at a given resolution, and
Forecast Primitive Compatibility (FPC) assesses whether a forecast primitive can
respond coherently to intervention, FAS assesses whether a slice should
participate in modeling and governance workflows at all. A slice may be
quantization-compatible and primitive-compatible in principle, yet still be
inadmissible for forecasting due to insufficient support or structurally hostile
error behavior.

Accordingly, FAS provides an upstream governance signal that answers three
practical questions:
\begin{enumerate}[label=(\roman*)]
  \item Is this slice sufficiently well-behaved to permit unrestricted modeling?
  \item If not, should forecasting be allowed only under constrained or
        conditional rules?
  \item In the most extreme cases, should this slice be excluded from forecasting
        and readiness control entirely?
\end{enumerate}

The remainder of this technical note formalizes the Forecast Admissibility
Surface. We define slice domains and baseline error anatomy, introduce a
deterministic admissibility taxonomy, and describe how FAS integrates into the
FRF without introducing new optimization targets, learned parameters, or
subjective tuning.
